\scanpage{183}\mspage{38}%
% archive sheet number: 183
% One continuous running-prose block; no visible paragraph indent or extra gap at Une fois.
in this first pre-derived-category approach, through knowledge of the fundamental group $\pi_1$, hence of the categories of twisted \unsure{constant|constant} coefficients $=\pi_1$-\unsure{module|modules}, the functors
% L5: the argument separator has a high point (about 344,446) and a separate lower comma-shaped stroke.
% Preferred semicolon; alternative comma with an unassigned detached point. Neither element is omitted.
$H^i(-\unsure{;|,}M)$ on these with \unsure{the|the} structure of \inserted{\unsure{cohomological|cohomological}} \unsure{functor|functors} \deleted{\unsure{ab|ab.}, of \unsure{functor|functors}} relative to exact sequences, \unsure{the|the} cup-product \deleted{\unsure{struc}}structure and so on --\raisebox{0pt}[0pt][0pt]{\parbox{0pt}{\figure{\begin{tikzpicture}[overlay,x=1em,y=1em]\fill (-0.55,-0.60) circle[radius=0.045em];\end{tikzpicture}}}} related by certain formal properties. Once one has at one's disposal the language of derived categories, the \inserted{subcategory of the derived category of} abelian complexes, \unsure{formed|form|formed|form-} \inserted{at the complexes of} sheaves with locally constant cohomology on $X$, with its triangulated structure and its \deleted{\unsure{m|sou|of}}multiplicative structure
% L13: detached compact point below tensor, approximately (1165,1373); origin and function unassigned.
$\overset{\mathbb{L}}{\otimes}\raisebox{0pt}[0pt][0pt]{\parbox{0pt}{\figure{\begin{tikzpicture}[overlay,x=1em,y=1em]\fill (-0.45,-0.95) ellipse[x radius=0.05em,y radius=0.07em];\end{tikzpicture}}}}$ (and possibly $\mathbb{R}\mathrm{Hom}$ \ldots) giving a more satisfactory candidate for hoping to recover the homotopy type. I have hardly any idea
% L17: separate short stroke below si (about 241,1836); two compact spots beside ils/permettent
% (about 802,1786 and 874,1810) have unassigned origin, irrespective of their color.
whether\raisebox{0pt}[0pt][0pt]{\parbox{0pt}{\figure{\begin{tikzpicture}[overlay,x=1em,y=1em]\draw[line width=0.6pt] (-0.45,-0.80) -- (-0.42,-0.87) -- (-0.45,-0.90) -- (-0.40,-0.97) -- (-0.37,-0.96);\end{tikzpicture}}}} or not they\raisebox{0pt}[0pt][0pt]{\parbox{0pt}{\figure{\begin{tikzpicture}[overlay,x=1em,y=1em]\fill (0.08,0.58) circle[radius=0.04em];\end{tikzpicture}}}} \raisebox{0pt}[0pt][0pt]{\parbox{0pt}{\figure{\begin{tikzpicture}[overlay,x=1em,y=1em]\fill (-0.12,0.05) ellipse[x radius=0.055em,y radius=0.07em];\end{tikzpicture}}}}allow one to recover it \inserted{in fact}, but on \raisebox{0pt}[0pt][0pt]{\parbox{0pt}{\figure{\begin{tikzpicture}[overlay,x=1em,y=1em]\fill (-0.10,-0.92) circle[radius=0.045em];\fill (0.06,-0.67) ellipse[x radius=0.045em,y radius=0.06em];\end{tikzpicture}}}}the other hand have hardly any doubt that by going all the way with the ``linearization'', that is to say by starting out in the
% L21 underline: first physical stroke about (318,2273)--(556,2261), second about (581,2262)--(673,2264).
% The first rule begins inside n: a short overlay preserves its terminal covered extent.
% The gap lies mainly under l; another short overlay preserves the second rule's beginning under l's terminal extent.
n\raisebox{0pt}[0pt][0pt]{\parbox{0pt}{\figure{\begin{tikzpicture}[overlay,x=1em,y=1em]\draw[line width=0.35pt] (-0.28,-0.16) -- (0,-0.16);\end{tikzpicture}}}}\uline{on-abe}l\raisebox{0pt}[0pt][0pt]{\parbox{0pt}{\figure{\begin{tikzpicture}[overlay,x=1em,y=1em]\draw[line width=0.35pt] (-0.10,-0.16) -- (0,-0.16);\end{tikzpicture}}}}\uline{ian} and working with the $(n+1)$-category (\unsure{without|its|san} \unsure{additional|additional} \unsure{structures|structure} on it!) of locally constant $n$-stacks\raisebox{0pt}[0pt][0pt]{\parbox{0pt}{\figure{\begin{tikzpicture}[overlay,x=1em,y=1em]\fill (0.08,0.62) circle[radius=0.04em];\end{tikzpicture}}}}
% L23: inner doubt concerns sur, not the clear X. The outer doubt concerns the insertion state.
% A curved guide attaches before pour; cancellation/older superposed ink remains a state alternative.
\unsure{\inserted{\unsure{on} $X$}} for all $n$
